% !TEX root = ../main.tex

\begin{abstract}
本项目实现面向真实网页任务的多模态 Agent 系统 \projectname{}。系统以自然语言任务和起始网页为输入，通过 Playwright 控制浏览器，组合原始截图、Set-of-Marks 标注图、DOM 候选元素、Accessibility Tree、页面文本及动作历史；模型采用 ReAct 风格单步决策，输出结构化 JSON 动作。执行侧完成动作校验、安全检查、进度监测及最终验证，并保存逐步 trace，支持 YAML 任务管理、批量运行、结果校验和 HTML 回放。

任务体系分为 \textbf{8 项核心评测任务与 1 项补充探索任务}。核心任务覆盖 CSRankings 信息检索、Bilibili 内容分析和广州旅游规划三个场景；CSRankings 跨方向排名与 faculty 统计单独作为探索性任务，其题目与判分设置仍需审查，不纳入核心完成率。DeepSeek V4 Pro 在每项核心任务上重复运行 10 次，共 80 次，记录到 76 次 harness 完成，完成率为 \textbf{95.00\%}。GPT-5.5 的视觉基线在 8 项核心任务上各运行一次，均完成，平均 9.75 步；其中两项 CSRankings 核心任务通过确定性校验，六项开放任务通过字段覆盖检查，仍需结合 trace 复核内容真实性。

项目另开展十种模型配置的横向实验，覆盖视觉与纯文本观察。原始批次包含核心和探索两部分，共 90 个计划组合；历史汇总记录为 66 个有效、22 个无效和 2 个未完成。由于历史汇总未完整区分每个模型的探索任务状态，这些数值仅用于说明原始批次覆盖，不作为八任务核心评测的分母或排行榜。

本报告给出系统架构、观察与动作空间、核心评测口径、探索任务范围及失败类型分析。完成率反映 Agent 主流程是否结束并生成答案；确定性任务采用答案规则校验，开放任务则需要字段、来源及轨迹证据共同复核。实验说明真实网页 Agent 的评估应同时报告任务范围、完成结果、校验方式与环境异常，避免将流程完成等同于答案正确。
\end{abstract}

\noindent\textbf{关键词：}网页 Agent；多模态大模型；Playwright；Set-of-Marks；ReAct；批量评测
